\documentclass[12pt,a4paper]{article}

\usepackage[utf8]{inputenc}
\usepackage[T1]{fontenc}
\usepackage[brazil]{babel}
\usepackage{geometry}
\usepackage{amsmath,amssymb}
\usepackage{booktabs}
\usepackage{longtable}
\usepackage{array}
\usepackage{enumitem}
\usepackage{hyperref}

\geometry{margin=2.5cm}
\setlist{itemsep=0.25em, topsep=0.4em}

\title{Modelos Matemáticos de Módulo de Resiliência Selecionados por Carvalho (2023)}
\author{}
\date{}

\begin{document}

\maketitle

\section*{Introdução}

Este texto apresenta, modelo a modelo, a formulação matemática dos quatorze modelos de comportamento do Módulo de Resiliência selecionados por Carvalho (2023), na dissertação intitulada \textit{Avaliação de modelos matemáticos de comportamento para o Módulo de Resiliência de Materiais de Pavimentação}. Os modelos foram reunidos pela autora no quadro-resumo da Tabela 5 da dissertação e utilizados para testes comparativos em bancos de dados de materiais de pavimentação.

O objetivo deste documento é organizar as expressões matemáticas, identificar o nome original de cada modelo e explicitar, de forma sintética, o significado das principais variáveis envolvidas.

\section*{Notação geral}

O Módulo de Resiliência é representado por:

\begin{equation}
MR
\end{equation}

em que $MR$ é o módulo de resiliência do material. As tensões básicas usadas nos modelos são:

\begin{equation}
\sigma_3
\end{equation}

\begin{equation}
\sigma_d
\end{equation}

em que $\sigma_3$ é a tensão confinante e $\sigma_d$ é a tensão desvio. A tensão desvio é definida por:

\begin{equation}
\sigma_d=\sigma_1-\sigma_3
\end{equation}

No ensaio triaxial convencional, admite-se:

\begin{equation}
\sigma_2=\sigma_3
\end{equation}

O primeiro invariante de tensões, ou soma das tensões principais, é dado por:

\begin{equation}
\theta=\sigma_1+\sigma_2+\sigma_3
\end{equation}

Como, no ensaio triaxial, $\sigma_2=\sigma_3$ e $\sigma_1=\sigma_d+\sigma_3$, resulta:

\begin{equation}
\theta=\sigma_d+3\sigma_3
\end{equation}

A tensão normal octaédrica é dada por:

\begin{equation}
\sigma_{oct}=\frac{\sigma_1+\sigma_2+\sigma_3}{3}=\frac{\theta}{3}
\end{equation}

A tensão cisalhante octaédrica é indicada, na formulação utilizada pela autora, por:

\begin{equation}
\tau_{oct}\approx0{,}471\sigma_d
\end{equation}

Além disso, utiliza-se:

\begin{equation}
P_a
\end{equation}

como pressão atmosférica de referência. Os parâmetros $k_1$, $k_2$ e $k_3$ são coeficientes experimentais ajustados por regressão.

\section{Dunlap (1963) -- modelo em função da tensão confinante}

Nome original apresentado pela autora: \textbf{Dunlap (1963)}.

\begin{equation}
MR=k_1\left(\frac{\sigma_3}{P_a}\right)^{k_2}
\end{equation}

Esse é um modelo de dois parâmetros, dependente apenas da tensão confinante $\sigma_3$. A formulação busca representar materiais cujo comportamento resiliente é controlado predominantemente pelo nível de confinamento. Em termos físicos, o modelo expressa a variação do módulo de resiliência com o aumento ou a redução da tensão lateral aplicada ao corpo de prova.

\section{Hicks (1970) -- modelo em função da tensão desvio}

Nome original apresentado pela autora: \textbf{Hicks (1970)}.

\begin{equation}
MR=k_1\sigma_d^{k_2}
\end{equation}

Esse modelo também possui dois parâmetros, mas considera apenas a tensão desvio $\sigma_d$ como variável independente. A autora associa esse tipo de formulação ao comportamento de solos finos ou coesivos, nos quais a tensão desvio pode exercer influência relevante sobre o módulo de resiliência.

\section{Witczak (1981)}

Nome original apresentado pela autora: \textbf{Witczak (1981)}.

Na forma apresentada no quadro-resumo da dissertação, a expressão é:

\begin{equation}
MR=k_1\frac{\theta^{k_2}}{P_a}\frac{\sigma_d^{k_3}}{P_a}
\end{equation}

Esse modelo introduz simultaneamente o primeiro invariante de tensões, $\theta$, e a tensão desvio, $\sigma_d$. Portanto, é mais completo do que modelos que utilizam somente $\sigma_3$ ou somente $\sigma_d$, pois considera tanto o efeito global do estado de tensões quanto o efeito desviador.

Convém observar que a forma escrita no quadro da dissertação pode gerar ambiguidade quanto à normalização por $P_a$. Uma forma normalizada usualmente empregada em formulações dessa família é:

\begin{equation}
MR=k_1\left(\frac{\theta}{P_a}\right)^{k_2}
\left(\frac{\sigma_d}{P_a}\right)^{k_3}
\end{equation}

Essa forma se aproxima da estrutura do modelo de Uzan.

\section{Uzan (1985)}

Nome original apresentado pela autora: \textbf{Uzan (1985)}.

\begin{equation}
MR=k_1\left(\frac{\theta}{P_a}\right)^{k_2}
\left(\frac{\sigma_d}{P_a}\right)^{k_3}
\end{equation}

É um modelo de três parâmetros, baseado no primeiro invariante de tensões e na tensão desvio. A autora observa que os modelos de Witczak e Uzan foram tratados como equivalentes em parte da análise, razão pela qual nem sempre aparecem separadamente nos resultados.

\section{Johnson et al. (1986) -- modelo do invariante de tensões}

Nome original apresentado pela autora: \textbf{Johnson et al. (1986)}.

\begin{equation}
MR=k_1\theta^{k_2}
\end{equation}

Esse modelo depende apenas do primeiro invariante de tensões $\theta$. Como:

\begin{equation}
\theta=\sigma_d+3\sigma_3
\end{equation}

no ensaio triaxial, mesmo sendo escrito como função de uma única variável, o modelo incorpora indiretamente os efeitos da tensão confinante e da tensão desvio.

\section{Witczak e Uzan (1988)}

Nome original apresentado pela autora: \textbf{Witczak e Uzan (1988)}.

\begin{equation}
MR=k_1P_a
\left(\frac{\theta}{P_a}\right)^{k_2}
\left(\frac{\tau_{oct}}{P_a}\right)^{k_3}
\end{equation}

Esse modelo substitui a tensão desvio direta por $\tau_{oct}$, isto é, pela tensão cisalhante octaédrica. Como:

\begin{equation}
\tau_{oct}\approx0{,}471\sigma_d
\end{equation}

continua havendo dependência do efeito desviador, porém por meio de uma grandeza derivada do estado tridimensional de tensões.

\section{Tam e Brown (1988)}

Nome original apresentado pela autora: \textbf{Tam e Brown (1988)}.

\begin{equation}
MR=k_1\left(\frac{\sigma_{oct}}{\sigma_d}\right)^{k_2}
\end{equation}

Como:

\begin{equation}
\sigma_{oct}=\frac{\theta}{3}
\end{equation}

pode-se escrever:

\begin{equation}
MR=k_1\left(\frac{\theta}{3\sigma_d}\right)^{k_2}
\end{equation}

Esse modelo difere dos anteriores porque usa uma razão entre a tensão normal octaédrica e a tensão desvio. No banco de dados de Ferreira (2008), analisado por Carvalho (2023), esse modelo apresentou o melhor desempenho geral, especialmente pelos critérios de $R^2$, $R^2$ ajustado e erro quadrático médio.

\section{Pezo (1993) -- modelo composto}

Nome original apresentado pela autora: \textbf{Pezo (1993)}.

\begin{equation}
MR=k_1P_a
\left(\frac{\sigma_3}{P_a}\right)^{k_2}
\left(\frac{\sigma_d}{P_a}\right)^{k_3}
\end{equation}

Esse é o chamado modelo composto, pois combina diretamente a tensão confinante $\sigma_3$ e a tensão desvio $\sigma_d$. No Brasil, esse tipo de formulação é amplamente utilizado em análises de módulo de resiliência. Na dissertação, o modelo de Pezo tem papel metodológico relevante, pois foi usado também para gerar valores de $MR$ em análises reversas de bancos de dados.

\section{Hopkins et al. (2001)}

Nome original apresentado pela autora: \textbf{Hopkins et al. (2001)}.

\begin{equation}
MR=k_1
\left(\frac{\sigma_3}{P_a}+1\right)^{k_2}
\left(\frac{\tau_{oct}}{P_a}+1\right)^{k_3}
\end{equation}

Esse modelo utiliza a tensão confinante normalizada e a tensão cisalhante octaédrica normalizada, ambas acrescidas de uma unidade. O acréscimo de $+1$ contribui para evitar problemas numéricos quando as tensões normalizadas assumem valores baixos.

\section{Ni et al. (2002)}

Nome original apresentado pela autora: \textbf{Ni et al. (2002)}.

\begin{equation}
MR=k_1P_a
\left(\frac{\sigma_3}{P_a}+1\right)^{k_2}
\left(\frac{\sigma_d}{P_a}+1\right)^{k_3}
\end{equation}

Esse modelo é semelhante ao modelo composto de Pezo, pois utiliza diretamente $\sigma_3$ e $\sigma_d$. A diferença é que ambos os termos normalizados são acrescidos de $+1$.

\section{NCHRP 1-28A (2004)}

Nome original apresentado no quadro da autora: \textbf{NCHRPI-28A (2004)}. Tecnicamente, a grafia mais adequada é \textbf{NCHRP 1-28A}.

\begin{equation}
MR=k_1P_a
\left(\frac{\theta}{P_a}\right)^{k_2}
\left(\frac{\sigma_d}{P_a}+1\right)^{k_3}
\end{equation}

Esse modelo combina o primeiro invariante de tensões, $\theta$, com a tensão desvio normalizada acrescida de uma unidade. Portanto, mistura uma variável global do estado de tensões com uma variável associada diretamente à solicitação desviadora.

\section{NCHRP 1-37A (2004)}

Nome original apresentado no quadro da autora: \textbf{NCHRPI-37A (2004)}. Tecnicamente, a grafia mais adequada é \textbf{NCHRP 1-37A}.

\begin{equation}
MR=k_1P_a
\left(\frac{\theta}{P_a}\right)^{k_2}
\left(\frac{\tau_{oct}}{P_a}+1\right)^{k_3}
\end{equation}

Esse modelo é semelhante ao NCHRP 1-28A, mas substitui a tensão desvio $\sigma_d$ pela tensão cisalhante octaédrica $\tau_{oct}$. Desse modo, o efeito desviador passa a ser representado por uma grandeza derivada do estado de tensões tridimensional.

\section{Ooi et al. (1) (2004)}

Nome original apresentado pela autora: \textbf{Ooi et al. (1) (2004)}.

\begin{equation}
MR=k_1P_a
\left(\frac{\theta}{P_a}+1\right)^{k_2}
\left(\frac{\sigma_d}{P_a}+1\right)^{k_3}
\end{equation}

Esse modelo utiliza $\theta$ e $\sigma_d$, ambos normalizados por $P_a$ e acrescidos de uma unidade. Em relação ao modelo NCHRP 1-28A, a diferença principal é que o termo associado ao invariante de tensões também recebe o acréscimo de $+1$.

\section{Ooi et al. (2) (2004)}

Nome original apresentado pela autora: \textbf{Ooi et al. (2) (2004)}.

\begin{equation}
MR=k_1P_a
\left(\frac{\theta}{P_a}\right)^{k_2}
\left(\frac{\tau_{oct}}{P_a}+1\right)^{k_3}
\end{equation}

Na forma apresentada pela autora, esse modelo é matematicamente igual ao modelo NCHRP 1-37A. Isso deve ser observado com atenção em uma revisão crítica, pois dois modelos aparecem com nomes diferentes, mas com a mesma estrutura formal no quadro-resumo da dissertação.

\section*{Quadro sintético dos quatorze modelos}

\renewcommand{\arraystretch}{1.35}
\begin{longtable}{p{4.3cm}p{10cm}}
\toprule
\textbf{Modelo} & \textbf{Expressão matemática} \\
\midrule
\endfirsthead
\toprule
\textbf{Modelo} & \textbf{Expressão matemática} \\
\midrule
\endhead
Dunlap (1963) & $MR=k_1\left(\dfrac{\sigma_3}{P_a}\right)^{k_2}$ \\
Hicks (1970) & $MR=k_1\sigma_d^{k_2}$ \\
Witczak (1981) & $MR=k_1\dfrac{\theta^{k_2}}{P_a}\dfrac{\sigma_d^{k_3}}{P_a}$ \\
Uzan (1985) & $MR=k_1\left(\dfrac{\theta}{P_a}\right)^{k_2}\left(\dfrac{\sigma_d}{P_a}\right)^{k_3}$ \\
Johnson et al. (1986) & $MR=k_1\theta^{k_2}$ \\
Witczak e Uzan (1988) & $MR=k_1P_a\left(\dfrac{\theta}{P_a}\right)^{k_2}\left(\dfrac{\tau_{oct}}{P_a}\right)^{k_3}$ \\
Tam e Brown (1988) & $MR=k_1\left(\dfrac{\sigma_{oct}}{\sigma_d}\right)^{k_2}$ \\
Pezo (1993) & $MR=k_1P_a\left(\dfrac{\sigma_3}{P_a}\right)^{k_2}\left(\dfrac{\sigma_d}{P_a}\right)^{k_3}$ \\
Hopkins et al. (2001) & $MR=k_1\left(\dfrac{\sigma_3}{P_a}+1\right)^{k_2}\left(\dfrac{\tau_{oct}}{P_a}+1\right)^{k_3}$ \\
Ni et al. (2002) & $MR=k_1P_a\left(\dfrac{\sigma_3}{P_a}+1\right)^{k_2}\left(\dfrac{\sigma_d}{P_a}+1\right)^{k_3}$ \\
NCHRP 1-28A (2004) & $MR=k_1P_a\left(\dfrac{\theta}{P_a}\right)^{k_2}\left(\dfrac{\sigma_d}{P_a}+1\right)^{k_3}$ \\
NCHRP 1-37A (2004) & $MR=k_1P_a\left(\dfrac{\theta}{P_a}\right)^{k_2}\left(\dfrac{\tau_{oct}}{P_a}+1\right)^{k_3}$ \\
Ooi et al. (1) (2004) & $MR=k_1P_a\left(\dfrac{\theta}{P_a}+1\right)^{k_2}\left(\dfrac{\sigma_d}{P_a}+1\right)^{k_3}$ \\
Ooi et al. (2) (2004) & $MR=k_1P_a\left(\dfrac{\theta}{P_a}\right)^{k_2}\left(\dfrac{\tau_{oct}}{P_a}+1\right)^{k_3}$ \\
\bottomrule
\end{longtable}

\section*{Observações críticas}

A lista dos quatorze modelos é coerente com o objetivo da dissertação, mas há algumas questões formais que merecem atenção em uma análise técnica mais rigorosa.

\begin{enumerate}
    \item A sigla aparece no quadro-resumo como \textbf{NCHRPI}, mas a forma tecnicamente adequada é \textbf{NCHRP}, referente ao \textit{National Cooperative Highway Research Program}.
    
    \item A autora indica quatorze modelos selecionados, mas, na metodologia, informa que foram efetivamente analisados treze, porque o modelo de Pezo foi usado em parte da geração dos dados e porque há sobreposição entre algumas formulações.
    
    \item O modelo de \textbf{Witczak (1981)} aparece no quadro com uma forma que pode gerar ambiguidade matemática quanto à posição dos expoentes em relação à normalização por $P_a$.
    
    \item Os modelos \textbf{NCHRP 1-37A (2004)} e \textbf{Ooi et al. (2) (2004)} aparecem, na forma apresentada pela autora, com a mesma expressão matemática. Isso deve ser tratado com cuidado, pois, se a equação é exatamente igual, a regressão tende a produzir o mesmo comportamento estrutural, salvo diferenças de implementação.
\end{enumerate}

\section*{Conclusão}

Os quatorze modelos selecionados por Carvalho (2023) podem ser organizados em famílias conforme a variável de tensão predominante: modelos dependentes de $\sigma_3$, modelos dependentes de $\sigma_d$, modelos baseados no invariante de tensões $\theta$, modelos que utilizam tensões octaédricas e modelos compostos que combinam mais de uma variável de tensão.

Em síntese, os modelos mais simples, como Dunlap, Hicks e Johnson et al., usam uma única variável de tensão. Os modelos mais completos, como Pezo, Uzan, Witczak e Uzan, NCHRP, Ni et al. e Ooi et al., incorporam simultaneamente diferentes componentes do estado de tensões. O modelo de Tam e Brown se distingue por trabalhar com uma razão entre tensão octaédrica e tensão desvio, tendo apresentado desempenho expressivo no banco de dados de Ferreira (2008), conforme a análise da autora.

\end{document}
